\chapter{Continuous Random Variables and Densities}

In the previous chapters, every calculation rested on a quiet assumption: the
variable was \emph{discrete}, taking values in a list you could write down.
The world, however, is full of quantities that vary smoothly: the time until a
radioactive atom decays, the height of a person, or the voltage on a circuit.
These quantities live on an interval of real numbers with no smallest step.

For a continuous variable, the probability of landing on any single exact value
is zero ($\P(X = x) = 0$). Instead of assigning probability to individual
points, we assign it to intervals as \emph{area} under a \emph{probability
density function} $f(x)$. The cumulative distribution function ties the
discrete and continuous views together, and expectation and variance become
integrals instead of sums.

\section{From bars to curves: the central picture}
\label{ch10:bars-to-curves}

Start where your intuition already is. Suppose you measure the waiting time
$X$ (in minutes) for a bus that is ``equally likely to come any time in the next
$10$ minutes.'' You cannot list the outcomes -- $X$ could be $3$, or $3.5$, or
$3.14159$, or anything in $[0,10]$. So let us approximate.

First, round $X$ to the nearest whole minute. Now you have a discrete variable on
$\set{0,1,\dots,10}$, and you can draw a familiar bar chart of probabilities. Each
bar has height roughly $1/10$. Now refine: round to the nearest half-minute. You
get twice as many bars, each half as tall, because each old minute splits into two
sub-intervals sharing its probability. Refine again to tenths of a minute: ten
times as many bars, each a tenth as tall. The bars are getting shorter and shorter
-- the \emph{height} $\P(X \in \text{tiny bin})$ shrinks toward zero -- but the
overall \emph{shape} is settling down to a flat plateau of constant height.

\begin{intuitionbox}
The probability in a single bin shrinks as the bins get finer, so heights are the
wrong thing to track. The right thing is height \emph{per unit width} -- the bar's
area divided by its width. That ratio does \emph{not} collapse to zero; it
converges to a stable curve. That stable curve is the density. Density is to a
continuous variable what the PMF was to a discrete one, but it lives on the
vertical axis as \emph{probability-per-length}, not probability.
\end{intuitionbox}

Figure~\ref{ch10:fig-refine} shows the refinement in action for a bell-shaped
example, where the limiting curve is more interesting than a flat line.

\begin{figure}[H]
\centering
\begin{tikzpicture}
\begin{axis}[
  width=0.92\textwidth, height=6.2cm,
  title style={font=\small},
  title={Histogram Bin Shrinking Toward Continuous Limiting Density $f(x)$},
  xlabel={$x$}, ylabel={density (area per unit $x$)},
  xmin=-3.6, xmax=3.6, ymin=0, ymax=0.5,
  domain=-3.5:3.5, samples=120,
  axis lines=left, ytick=\empty,
  legend style={at={(0.02,0.98)},anchor=north west,draw=none,fill=none,font=\small},
]
% coarse histogram bars (width 1)
\addplot+[ybar interval, mark=no, fill=histfill!45, draw=histfill!80!black, opacity=0.55,
  forget plot]
  coordinates {(-3.5,0.0060)(-2.5,0.0606)(-1.5,0.2417)(-0.5,0.3829)(0.5,0.2417)
               (1.5,0.0606)(2.5,0.0060)(3.5,0)};
% the limiting density: standard normal
\addplot[densityred, very thick, smooth] {1/sqrt(2*pi)*exp(-x^2/2)};
\addlegendentry{limiting density $f(x)$}
\end{axis}
\end{tikzpicture}
\caption{Coarse probability bars (width $1$) with their heights equal to the
probability in each bin. As bins shrink, the bars hug a smooth curve when you plot
\emph{probability per unit length}. That curve is the probability density.}
\label{ch10:fig-refine}
\end{figure}

The picture suggests the whole plan. A density $f$ is a curve such that the
\emph{area} under it over an interval equals the probability of landing in that
interval:
\[
  \P(a \le X \le b) \;=\; \text{area under } f \text{ from } a \text{ to } b
  \;=\; \int_a^b f(x)\,\dd x .
\]
Probabilities are areas -- that one equation carries most of continuous
probability. Figure~\ref{ch10:fig-area} shows it directly: the probability of
landing in $[a,b]$ is the shaded region under the curve.

\begin{figure}[H]
\centering
\begin{tikzpicture}
\begin{axis}[
  width=0.85\textwidth, height=5.4cm,
  title style={font=\small},
  title={Probability as Area Under the Density Curve: $\P(a \le X \le b) = \int_a^b f(x)\,\dd x$},
  xlabel={$x$}, ylabel={$f(x)$},
  xmin=-3.4, xmax=3.4, ymin=0, ymax=0.48,
  domain=-3.3:3.3, samples=120,
  axis lines=left, ytick=\empty,
  xtick={-1,1.4}, xticklabels={$a$,$b$},
]
\addplot[name path=curve, densityred, very thick] {1/sqrt(2*pi)*exp(-x^2/2)};
\path[name path=axis] (axis cs:-1,0) -- (axis cs:1.4,0);
\addplot[densityred!18] fill between[of=curve and axis, soft clip={domain=-1:1.4}];
\node at (axis cs:0.2,0.13) {\small $\P(a\le X\le b)$};
\end{axis}
\end{tikzpicture}
\caption{Probability as area. For a continuous variable the chance of landing in
$[a,b]$ is the shaded area under the density $f$ between $a$ and $b$, i.e.
$\int_a^b f(x)\,\dd x$. Total area under the whole curve is $1$.}
\label{ch10:fig-area}
\end{figure}

\section{Why \texorpdfstring{$\P(X=x)=0$}{P(X=x)=0}, and why that is fine}
\label{ch10:point-mass-zero}

Here is the fact that trips up everyone the first time. If $X$ is continuous, then
for \emph{every} particular value $x$,
\[
  \P(X = x) = 0 .
\]
Not ``very small.'' Exactly zero. Yet $X$ does take \emph{some} value every time
you sample it. How can every outcome have probability zero, while one of them
always happens?

\begin{happeningbox}
The trap is importing intuition from the discrete world, where ``probability zero''
means ``impossible.'' In the continuum the two notions come apart. Probability is
\emph{area}, and a single point is a line segment of width zero, so it bounds zero
area no matter how tall the density is there. Think of probability as a
\emph{liquid} of total volume $1$ poured over the real line. Ask how much liquid
sits exactly over the single point $x=0.5$. None -- a point has no width to hold any.
But the liquid over the interval $[0.49,0.51]$ is a real, positive amount. In the
continuum, \emph{probability lives in intervals, not in points.} ``Probability $0$''
no longer means ``cannot happen''; it means ``occupies no length.''
\end{happeningbox}

We can prove $\P(X=x)=0$ directly from the area picture, and the proof is
illuminating. Pin down a value $x$ and trap it in a shrinking interval. For any
$\eps > 0$,
\[
  \P(X = x) \;\le\; \P\!\left(x - \tfrac{\eps}{2} \le X \le x + \tfrac{\eps}{2}\right)
  \;=\; \int_{x-\eps/2}^{x+\eps/2} f(t)\,\dd t .
\]
If $f$ is bounded near $x$, say $f \le M$ there, this integral is at most $M\eps$,
which we can drive to $0$ by shrinking $\eps$. A nonnegative number smaller than
every positive number must be $0$. So $\P(X=x)=0$.

\begin{trapbox}
Because individual points carry no probability, the difference between $<$ and
$\le$ \emph{vanishes} for continuous variables:
\[
  \P(a \le X \le b) = \P(a < X < b) = \P(a \le X < b) = \P(a < X \le b).
\]
Including or excluding an endpoint changes the answer by $\P(X=a)$ or $\P(X=b)$,
which are both $0$. This is a genuine convenience -- but it is \emph{false} for
discrete variables, where endpoints can carry real mass. Always know which world
you are in.
\end{trapbox}

\section{The probability density function, precisely}
\label{ch10:pdf-def}

We can now state the definition the picture has been pointing at.

\begin{definition}[Continuous random variable and density]
\label{ch10:def-pdf}
A random variable $X$ is \emph{continuous} if there is a function
$f:\R \to [0,\infty)$, called its \emph{probability density function} (pdf), such
that for every interval $[a,b]$,
\[
  \P(a \le X \le b) = \int_a^b f(x)\,\dd x .
\]
A pdf must satisfy two requirements, the continuous analogues of the PMF's rules:
\begin{enumerate}
\item \textbf{Nonnegativity:} $f(x) \ge 0$ for all $x$ (areas of probability cannot
      be negative);
\item \textbf{Normalization:} $\displaystyle \int_{-\infty}^{\infty} f(x)\,\dd x = 1$
      (all the probability is somewhere).
\end{enumerate}
\end{definition}

\begin{trapbox}
A density is \emph{not} a probability, and it is not capped at $1$. The value
$f(x)$ has units of probability \emph{per unit $x$}. It can be huge -- a density
sharply concentrated on $[0,0.01]$ might reach a height of $100$ -- as long as the
total area stays $1$. Never say ``$f(x)$ is the probability that $X=x$.'' The
probability of any exact value is $0$. The right reading is local: $f(x)\,\dd x$ is
approximately the probability that $X$ lands in a tiny window of width $\dd x$
around $x$.
\end{trapbox}

That last sentence is worth promoting to a working tool, because it is how
practitioners actually think.

\begin{focusbox}
\textbf{The infinitesimal reading of a density.} For a small width $\dd x$,
\[
  \P\!\left(x \le X \le x + \dd x\right) \;\approx\; f(x)\,\dd x .
\]
Read $f(x)$ as ``how densely probability is packed near $x$.'' Where $f$ is tall,
outcomes cluster; where $f$ is short, they are rare; where $f=0$, they never occur.
Ratios are meaningful even though absolute values are not: if $f(2) = 3 f(7)$, then
a tiny window near $2$ is three times as likely as an equally tiny window near $7$.
\end{focusbox}

\begin{walkbox}
\textbf{Pause and think: Can a density exceed $1$?} Suppose a random variable
$X$ has constant density on the narrow interval $[0, 0.2]$ and is $0$
everywhere else. What must the value of $f(x)$ be on that interval?

\emph{Answer:} Normalization demands that the total area under $f$ equals $1$.
Here, the region is a rectangle of width $0.2$, so $\text{area} = \text{height}
\times 0.2 = 1$, giving $\text{height} = 1/0.2 = 5$. The density is $5$ -- five
times higher than $1$! Yet the probability of any sub-interval, such as
$[0, 0.1]$, is $\text{area} = 5 \times 0.1 = 0.5 \le 1$. Densities are
\emph{rates} (probability per unit width), not probabilities. As long as the
base is narrow, the rate can reach any height.
\end{walkbox}

\begin{example}[Checking and using a density]
\label{ch10:ex-quadratic}
Suppose $X$ has the candidate density
\[
  f(x) = \begin{cases} c\,x^2, & 0 \le x \le 1,\\[2pt] 0, & \text{otherwise.}\end{cases}
\]
First find $c$ so that $f$ is a legitimate density, then compute
$\P(X \le \tfrac12)$.

\textbf{Normalize.} We need the total area to be $1$:
\[
  1 = \int_{-\infty}^{\infty} f(x)\,\dd x = \int_0^1 c\,x^2 \,\dd x
    = c \cdot \frac{x^3}{3}\Big|_0^1 = \frac{c}{3}
  \quad\Longrightarrow\quad c = 3 .
\]
So $f(x) = 3x^2$ on $[0,1]$. It is nonnegative there, so it is a valid density.

\textbf{Compute the probability.} Probability is area:
\[
  \P\!\left(X \le \tfrac12\right) = \int_0^{1/2} 3x^2 \,\dd x
    = x^3 \Big|_0^{1/2} = \left(\tfrac12\right)^3 = \tfrac18 .
\]
Notice the answer is much less than $1/2$, even though $1/2$ is the midpoint of the
range. The density $3x^2$ grows toward $x=1$, so probability piles up on the right.
\end{example}

\section{The cumulative distribution function: the bridge}
\label{ch10:cdf}

There is one object that works for \emph{every} random variable -- discrete,
continuous, or a mix -- and ties the whole subject together. It is the
\emph{cumulative distribution function} (CDF).

\begin{definition}[CDF]
\label{ch10:def-cdf}
The CDF of any random variable $X$ is
\[
  F(x) = \P(X \le x), \qquad x \in \R .
\]
It accumulates all the probability lying at or below $x$.
\end{definition}

For a continuous $X$, accumulating probability means accumulating area, so the CDF
is the running integral of the density:
\[
  F(x) = \P(X \le x) = \int_{-\infty}^{x} f(t)\,\dd t .
\]
Every CDF, regardless of type, obeys three structural rules that come straight from
the axioms of probability:
\begin{enumerate}
\item $F$ is \textbf{nondecreasing} (more territory, more accumulated probability);
\item $\displaystyle \lim_{x \to -\infty} F(x) = 0$ and
      $\displaystyle \lim_{x \to +\infty} F(x) = 1$ (eventually we have collected
      none, then all, of the probability);
\item $F$ is \textbf{right-continuous} (approaching any point from the right
      settles onto the value at that point, $\lim_{h \to 0^+} F(x+h) = F(x)$,
      because $F(x)$ includes the exact boundary point $x$).
\end{enumerate}
The defining contrast between the two worlds is visual. A discrete CDF is a
\emph{staircase}: flat stretches punctuated by vertical jumps, and the height of
the jump at $x$ is exactly $\P(X=x)$. A continuous CDF has \emph{no jumps} -- it is
a smooth ramp -- precisely because every point mass is $0$. Figure~\ref{ch10:fig-cdf}
puts the two side by side.

\begin{figure}[H]
\centering
\begin{tikzpicture}
\begin{axis}[
  name=disc, width=0.46\textwidth, height=5.2cm,
  title style={font=\small},
  title={Discrete CDF (Staircase Jumps)},
  xlabel={$x$}, ylabel={$F(x)$},
  xmin=-0.5, xmax=4.5, ymin=0, ymax=1.08,
  axis lines=left, ytick={0,0.25,0.5,0.75,1},
]
% staircase for X in {0,1,2,3} each prob 1/4
\addplot[distblue, very thick] coordinates {(-0.5,0)(0,0)};
\addplot[distblue, very thick] coordinates {(0,0.25)(1,0.25)};
\addplot[distblue, very thick] coordinates {(1,0.5)(2,0.5)};
\addplot[distblue, very thick] coordinates {(2,0.75)(3,0.75)};
\addplot[distblue, very thick] coordinates {(3,1)(4.5,1)};
\addplot[only marks, mark=*, distblue] coordinates {(0,0.25)(1,0.5)(2,0.75)(3,1)};
\addplot[only marks, mark=o, distblue] coordinates {(0,0)(1,0.25)(2,0.5)(3,0.75)};
\end{axis}
\begin{axis}[
  at={(disc.south east)}, anchor=south west, xshift=1.2cm,
  width=0.46\textwidth, height=5.2cm,
  title style={font=\small},
  title={Continuous CDF (Smooth Ramp)},
  xlabel={$x$}, ylabel={$F(x)$},
  xmin=-0.2, xmax=1.2, ymin=0, ymax=1.08,
  axis lines=left, ytick={0,0.25,0.5,0.75,1},
  domain=0:1, samples=80,
]
% CDF of f=3x^2 on [0,1] is x^3
\addplot[densityred, very thick] {x^3};
\addplot[densityred, very thick] coordinates {(-0.2,0)(0,0)};
\addplot[densityred, very thick] coordinates {(1,1)(1.2,1)};
\end{axis}
\end{tikzpicture}
\caption{Left: a discrete CDF is a staircase, with jump heights equal to point
masses. Right: a continuous CDF (here $F(x)=x^3$ for the density $f=3x^2$) is a
jumpless ramp -- every point carries probability $0$, so there is nothing to jump.}
\label{ch10:fig-cdf}
\end{figure}

\begin{walkbox}
\textbf{Pause and think: Reading probabilities directly from the CDF.} A continuous
variable $X$ has CDF $F$. You are given $F(3) = 0.40$ and $F(5) = 0.85$. What is
$\P(3 < X \le 5)$? Does it matter whether the endpoints are included or excluded?

\emph{Answer:} Since $F(5)$ accounts for all probability up to $5$ and $F(3)$ accounts
for everything up to $3$, the slice in between is $F(5) - F(3) = 0.85 - 0.40 = 0.45$.
Because $X$ is continuous, individual points carry zero probability, so $\P(3 \le X
\le 5) = \P(3 < X < 5) = 0.45$ as well. You obtained the exact probability without
integrating, because the CDF already completed the accumulation.
\end{walkbox}

\subsection{Getting the density back from the CDF}
\label{ch10:density-from-cdf}

The CDF integrates the density. The Fundamental Theorem of Calculus runs that
machine in reverse: \emph{differentiate the CDF to recover the density.}

\begin{theorem}[Density is the derivative of the CDF]
\label{ch10:thm-fdash}
If $X$ has CDF $F$ that is differentiable at $x$, then its density is
\[
  f(x) = F'(x) .
\]
\end{theorem}

\begin{whybox}
This is a direct application of the Fundamental Theorem of Calculus in a
probabilistic setting. Start from the definition of the derivative of $F$ and read
each piece probabilistically:
\[
  \begin{aligned}
  F'(x) &= \lim_{h \to 0} \frac{F(x+h) - F(x)}{h}
         = \lim_{h \to 0} \frac{\P(x < X \le x + h)}{h}\\[2pt]
        &= \lim_{h \to 0} \frac{\text{probability in a window of width } h}{h}.
  \end{aligned}
\]
That final ratio is exactly \emph{probability per unit length} near $x$ -- which is
the very thing we have been calling density. So $F'(x)=f(x)$ says: the slope of the
accumulated-probability curve at $x$ \emph{is} the local density. Steep ramp $=$
probability arriving fast $=$ tall density. Flat ramp $=$ no probability here $=$
density $0$. The CDF and pdf are two views of one object, related by the same
integrate/differentiate duality you met for area and slope in calculus.
\end{whybox}

\begin{connectionbox}
This is the cleanest place in all of probability to \emph{see} the Fundamental
Theorem of Calculus at work. $F = \int f$ and $f = F'$ are an
integrate--differentiate pair: the CDF is the antiderivative (the ``area-so-far''
function) of the density, and the density is the slope of the CDF. Everything you
know about antiderivatives -- that they are determined up to a constant, that
definite integrals are computed by evaluating an antiderivative at the endpoints --
transfers directly. In particular
\[
  \P(a \le X \le b) = \int_a^b f = F(b) - F(a),
\]
which is just ``evaluate the antiderivative at the endpoints.'' See the companion
calculus volume for the FTC; here it is the engine that converts between the two
faces of a distribution.
\end{connectionbox}

\begin{example}[Round trip: pdf $\to$ CDF $\to$ pdf]
\label{ch10:ex-roundtrip}
Take again $f(x) = 3x^2$ on $[0,1]$. Build the CDF by integrating:
\[
  F(x) = \int_0^x 3t^2 \,\dd t = x^3 \quad \text{for } 0 \le x \le 1,
\]
with $F(x)=0$ for $x<0$ and $F(x)=1$ for $x>1$. Check the endpoints: $F(0)=0$,
$F(1)=1$, good. Now differentiate to recover the density: for $0<x<1$,
$F'(x) = 3x^2 = f(x)$. The loop closes. As a bonus, the CDF makes
Example~\ref{ch10:ex-quadratic} instant: $\P(X\le \tfrac12) = F(\tfrac12)
= (\tfrac12)^3 = \tfrac18$, no integral required at the point of use.
\end{example}

\begin{walkbox}
You are told $X$ has CDF $F(x) = 1 - e^{-2x}$ for $x \ge 0$ (and $0$ for $x<0$).
What is the density? What is $\P(X > 1)$?

\emph{Density:} differentiate. $f(x) = F'(x) = 2e^{-2x}$ for $x \ge 0$. (This is the
$\Expo(2)$ density -- the exponential distribution of a later chapter.) \emph{Tail
probability:} the CDF gives it directly, $\P(X>1) = 1 - F(1) = 1 - (1 - e^{-2})
= e^{-2} \approx 0.135$. Carrying the CDF around is often faster than re-integrating
the density.
\end{walkbox}

\section{Expectation and variance as integrals}
\label{ch10:expectation}

For a discrete variable you computed the mean by summing value $\times$
probability:
$\E[X] = \sum_x x\,\P(X=x)$. In the continuum the point masses $\P(X=x)$ are all
$0$, so that sum is hopeless. But recall the infinitesimal reading: the probability
in a sliver of width $\dd x$ near $x$ is $f(x)\,\dd x$. Replace the discrete mass
$\P(X=x)$ by the sliver mass $f(x)\,\dd x$, and replace the sum by an integral. The
recipe ``value times probability, then add up'' becomes:

\begin{definition}[Expectation and variance of a continuous variable]
\label{ch10:def-EV}
For a continuous $X$ with density $f$,
\[
  \E[X] = \int_{-\infty}^{\infty} x\,f(x)\,\dd x,
  \qquad
  \E[g(X)] = \int_{-\infty}^{\infty} g(x)\,f(x)\,\dd x,
\]
the latter being the continuous Law of the Unconscious Statistician (LOTUS) -- the
rule that allows you to calculate the expected value of a transformed quantity
$g(X)$ directly against the original density $f(x)$, without having to derive the
distribution of the new variable $g(X)$ first. The variance is the mean squared
deviation, exactly as before:
\[
  \Var(X) = \E\!\big[(X-\mu)^2\big] = \int_{-\infty}^{\infty} (x-\mu)^2 f(x)\,\dd x
          = \E[X^2] - \big(\E[X]\big)^2,
\]
where $\mu = \E[X]$, and $\SD(X) = \sqrt{\Var(X)}$.
\end{definition}

\begin{connectionbox}
The transition $\sum \rightsquigarrow \int$ is not an analogy -- it is the literal
definition of the integral. The Riemann integral $\int x f(x)\,\dd x$ \emph{is} the
limit of sums $\sum_i x_i\, f(x_i)\,\Delta x$ over finer and finer partitions. Each
term $x_i \cdot \big(f(x_i)\Delta x\big)$ is ``a value times the probability of a
bin,'' precisely the discrete expectation of the binned-up variable. Letting the
bins shrink turns the discrete mean into the integral. So continuous expectation is
exactly discrete expectation in the limit of vanishing bin width. (The fully
general definition, valid even when no density exists, uses the Lebesgue integral
against the distribution -- which slices the vertical axis of probabilities rather
than chopping the horizontal domain into tiny intervals -- developed in advanced
measure theory.)
\end{connectionbox}

\begin{intuitionbox}
The mean $\E[X]$ is the \emph{balance point} of the density, the place where a
cardboard cutout of the region under $f$ would balance on a knife edge. In physics,
a lever balances when the net torque about the fulcrum is zero. Here, the lever arm
from the fulcrum $\mu$ to position $x$ is $(x - \mu)$, and the mass element is
$f(x)\,\dd x$. The net torque about $\mu$ is
\[
  \int_{-\infty}^{\infty} (x - \mu) f(x)\,\dd x
  \;=\; \int_{-\infty}^{\infty} x f(x)\,\dd x - \mu \int_{-\infty}^{\infty} f(x)\,\dd x
  \;=\; \E[X] - \mu \cdot 1 \;=\; 0 .
\]
The mean $\mu = \E[X]$ is literally the unique point where total torque vanishes.
The variance is its \emph{rotational inertia} (moment of inertia) about that balance
point: mass far from center contributes as the square of its distance, so $\Var$
measures spread, heavily weighting the tails. These mechanical pictures hold
identically in the discrete and continuous cases; only the bookkeeping (sum vs.\
integral) changes.
\end{intuitionbox}

\begin{example}[Mean and variance of $f=3x^2$]
\label{ch10:ex-EV}
With $f(x)=3x^2$ on $[0,1]$:
\[
  \E[X] = \int_0^1 x\cdot 3x^2 \,\dd x = \int_0^1 3x^3 \,\dd x
        = \frac{3x^4}{4}\Big|_0^1 = \frac34 .
\]
This sits to the right of the midpoint $\tfrac12$, confirming that the rising
density tips the balance point toward $1$, as illustrated in Figure~\ref{ch10:fig-balance}.
For the variance, get the second moment first:
\[
  \E[X^2] = \int_0^1 x^2 \cdot 3x^2 \,\dd x = \int_0^1 3x^4 \,\dd x
          = \frac{3x^5}{5}\Big|_0^1 = \frac35 .
\]
Then
\[
  \Var(X) = \E[X^2] - \big(\E[X]\big)^2 = \frac35 - \left(\frac34\right)^2
          = \frac35 - \frac{9}{16} = \frac{48 - 45}{80} = \frac{3}{80}
          = 0.0375,
\]
so $\SD(X) = \sqrt{3/80} \approx 0.194$.
\end{example}

\begin{figure}[H]
\centering
\begin{tikzpicture}
\begin{axis}[
  width=0.85\textwidth, height=5.4cm,
  title style={font=\small},
  title={Expectation as Physical Balance Point (Center of Mass): $\int (x-\mu)f(x)\,\dd x = 0$},
  xlabel={$x$}, ylabel={$f(x) = 3x^2$},
  xmin=-0.08, xmax=1.08, ymin=0, ymax=3.4,
  domain=0:1, samples=100,
  axis lines=left, ytick=\empty,
  xtick={0, 0.5, 0.75, 1}, xticklabels={$0$, $0.5$, $\mu = 0.75$, $1$},
]
\addplot[name path=curve, sourceblue, very thick] {3*x^2};
\path[name path=axis] (axis cs:0,0) -- (axis cs:1,0);
\addplot[sourceblue!20] fill between[of=curve and axis, soft clip={domain=0:1}];
\addplot[black, dashed, thick] coordinates {(0.75,0)(0.75,1.6875)};
\addplot[only marks, mark=*, mark size=2.5pt, sourceblue!70!black] coordinates {(0.75,1.6875)};
\node[above=3pt] at (axis cs:0.75, 1.6875) {\small $(\mu, f(\mu))$};
\node[align=center, font=\small] at (axis cs:0.35, 1.8) {long lever arm\\($x < \mu$)};
\draw[->, semithick, gray!70] (axis cs:0.35, 1.4) -- (axis cs:0.35, 0.5);
\node[align=center, font=\scriptsize] at (axis cs:0.91, 0.7) {short arm\\($x > \mu$)};
\end{axis}
\end{tikzpicture}
\caption{The expectation $\E[X] = 0.75$ is the balance point (center of mass) of the
density region. The light left tail sits farther from $\mu$ with a longer lever arm;
the heavy right peak sits closer to $\mu$ with a shorter lever arm. The net torque
$\int_0^1 (x-\mu)f(x)\,\dd x$ about $\mu = 0.75$ balances exactly to zero.}
\label{ch10:fig-balance}
\end{figure}

\begin{aestheticbox}
Every theorem you proved about expectation in the discrete chapters -- linearity
$\E[aX+b] = a\E[X]+b$, the variance scaling $\Var(aX+b) = a^2\Var(X)$, the identity
$\Var(X)=\E[X^2]-\E[X]^2$ -- holds \emph{verbatim} for continuous variables. Their
proofs used only the linearity of the operation that combines values with weights,
and integration is linear exactly as summation is. That is the payoff of a unified
theory: you do not relearn expectation, you reuse it, swapping $\sum$ for $\int$.
\end{aestheticbox}

\section{The uniform distribution: the base case}
\label{ch10:uniform}

Every subject has a simplest nontrivial example worth knowing cold. For continuous
probability it is the \emph{uniform} distribution, the formal version of ``a number
chosen completely at random from an interval, with no part favored.''

\begin{definition}[Uniform distribution]
\label{ch10:def-unif}
$X \sim \Unif(a,b)$ has constant density on $[a,b]$:
\[
  f(x) = \begin{cases} \dfrac{1}{b-a}, & a \le x \le b,\\[6pt] 0, & \text{otherwise.}\end{cases}
\]
The height $1/(b-a)$ is forced by normalization: a rectangle of width $b-a$ and area
$1$ must have that height.
\end{definition}

Its CDF is a straight ramp -- probability accumulates at a constant rate:
\[
  F(x) = \begin{cases} 0, & x < a, \\[2pt] \dfrac{x-a}{b-a}, & a \le x \le b, \\[6pt] 1, & x > b. \end{cases}
\]
Figure~\ref{ch10:fig-unif} shows the flat density and its ramp CDF together.

\begin{figure}[H]
\centering
\begin{tikzpicture}
\begin{axis}[
  name=u1, width=0.46\textwidth, height=4.6cm,
  title style={font=\small},
  title={$\Unif(0,1)$ Density $f(x)$},
  xlabel={$x$}, ylabel={$f(x)$},
  xmin=-0.4, xmax=1.4, ymin=0, ymax=1.4,
  axis lines=left,
]
\addplot[fill=densityred!18, draw=none] coordinates {(0,0)(0,1)(1,1)(1,0)(0,0)};
\addplot[densityred, very thick] coordinates {(-0.4,0)(0,0)};
\addplot[densityred, very thick] coordinates {(0,1)(1,1)};
\addplot[densityred, very thick] coordinates {(1,0)(1.4,0)};
\addplot[densityred, very thick, dashed] coordinates {(0,0)(0,1)};
\addplot[densityred, very thick, dashed] coordinates {(1,0)(1,1)};
\end{axis}
\begin{axis}[
  at={(u1.south east)}, anchor=south west, xshift=1.2cm,
  width=0.46\textwidth, height=4.6cm,
  title style={font=\small},
  title={$\Unif(0,1)$ CDF $F(x)$},
  xlabel={$x$}, ylabel={$F(x)$},
  xmin=-0.4, xmax=1.4, ymin=0, ymax=1.1,
  axis lines=left,
]
\addplot[distblue, very thick] coordinates {(-0.4,0)(0,0)};
\addplot[distblue, very thick] coordinates {(0,0)(1,1)};
\addplot[distblue, very thick] coordinates {(1,1)(1.4,1)};
\end{axis}
\end{tikzpicture}
\caption{The uniform $\Unif(0,1)$: a flat density (left) of constant height $1$, and
its CDF (right), a straight ramp from $0$ to $1$. Constant density means probability
accumulates at a constant rate.}
\label{ch10:fig-unif}
\end{figure}

\begin{example}[Mean and variance of $\Unif(a,b)$]
\label{ch10:ex-unif-EV}
By symmetry the balance point is the midpoint, $\E[X] = (a+b)/2$, but let us see the
integral confirm it:
\[
  \E[X] = \int_a^b x\,\frac{1}{b-a}\,\dd x
        = \frac{1}{b-a}\cdot \frac{x^2}{2}\Big|_a^b
        = \frac{b^2-a^2}{2(b-a)} = \frac{a+b}{2}.
\]
For the variance, first get the second moment,
\[
  \E[X^2] = \int_a^b x^2\,\frac{1}{b-a}\,\dd x
          = \frac{1}{b-a}\cdot \frac{x^3}{3}\Big|_a^b
          = \frac{b^3-a^3}{3(b-a)} = \frac{a^2+ab+b^2}{3},
\]
using $b^3-a^3 = (b-a)(a^2+ab+b^2)$. Then subtract the squared mean:
\[
  \Var(X) = \E[X^2] - \big(\E[X]\big)^2
          = \frac{a^2+ab+b^2}{3} - \frac{(a+b)^2}{4}
          = \frac{4(a^2+ab+b^2) - 3(a+b)^2}{12}
          = \frac{(b-a)^2}{12}.
\]
The factor $1/12$ is worth remembering -- it is a recurring constant in noise
modeling and in the analysis of rounding error.
\end{example}

\begin{applicationbox}
\textbf{Quantization noise and the $1/12$ rule.} When an analog signal is rounded
to the nearest digital level spaced $\Delta$ apart, the rounding error is, to good
approximation, $\Unif(-\Delta/2,\,\Delta/2)$. By Example~\ref{ch10:ex-unif-EV} its
mean is $0$ and its variance is $\Delta^2/12$. That single fact is the foundation of
the signal-processing rule that each extra bit of resolution (halving $\Delta$)
improves the signal-to-noise ratio by about $6$~dB: halving $\Delta$ divides the
noise variance by $4$, i.e.\ the noise standard deviation by $2$, and
$20\log_{10}2 \approx 6$. A toy model of rounding, made precise, drives real
engineering decisions about how many bits a converter needs.
\end{applicationbox}

\begin{historybox}
The idea that probability could be \emph{geometric} -- an area or a length rather
than a count -- predates the formal theory of densities by two centuries. When
Pierre de Fermat, Blaise Pascal, and Christiaan Huygens founded probability in the
1650s, they worked exclusively with discrete games of chance: rolling dice, drawing
cards, and tallying finite combinations. But natural scientists and astronomers
faced a conceptual wall: how do you assign a probability to a continuous quantity
like an angle, an arrival time, or a spatial coordinate? On a continuum, counting
outcomes collapses completely, because every single point carries probability zero.
If every individual outcome has zero chance, how does the collection of all outcomes
sum to one?

The earliest probabilists had only classical Euclidean geometry and newborn
calculus -- tools built to compute physical lengths and areas, but lacking any
formal rules for what kinds of continuous sets could legitimately carry probability.
In 1733, Georges-Louis Leclerc, Comte de Buffon, posed a breakthrough puzzle:
drop a needle of length $\ell$ onto a floor made of parallel wooden planks of width
$d$ (with $\ell \le d$); what is the chance the needle crosses a crack? Buffon
realized that the needle's state could be described by two continuous quantities --
the distance from its center to the nearest crack and the angle it forms with the
lines. By plotting this two-dimensional region of outcomes, he showed that the
probability equals the ratio of favorable area to total area, yielding $2\ell/(\pi d)$.
Buffon's needle was the first continuous probability problem ever solved, and it
even offered an experimental way to approximate $\pi$ by tossing needles.

Yet treating probability simply as geometric area soon raised controversies.
In 1889, Joseph Bertrand showed that the innocent phrase ``choose a chord of a circle
at random'' leads to three conflicting answers ($1/2$, $1/3$, or $1/4$) depending on
which geometric coordinates you choose to make uniform. The paradox proved that
continuous probability cannot rely on unguided geometric intuition; one must
explicitly specify the probability density.

The crisis was finally resolved in 1933, when Andrey Kolmogorov published his
axiomatic foundation of probability based on measure theory -- the branch of
analysis that formalizes the notion of size (lengths, areas, and volumes) for
abstract sets. By defining probability as a normalized measure that assigns a volume
between $0$ and $1$ to valid subsets (events), Kolmogorov unified discrete sums and
continuous integrals into a single, rigorous framework. This provided
the mathematical foundation for twentieth-century physics, including quantum mechanics
and statistical mechanics.

Even today, subtle foundational questions remain open at the boundary of continuous
conditioning. The Borel--Kolmogorov paradox reveals that conditioning on an event of
probability zero (such as restricting a random point on a sphere to lie on a great
circle) depends on the coordinate system chosen to define the limit. Resolving how
to condition on null sets in a coordinate-free way remains an active area of
research in Bayesian inference, causal modeling, and geometric probability.
\end{historybox}

\section{Quantiles: inverting the CDF}
\label{ch10:quantiles}

The CDF answers ``given a value $x$, what fraction of probability lies below it?''
Quantiles answer the reverse question, which is often the one you actually want:
``given a fraction $p$, what value has that fraction of probability below it?'' The
$90$th percentile of a test, the median income, the $0.975$ critical value behind a
confidence interval -- all are quantiles.

\begin{definition}[Quantile function]
\label{ch10:def-quantile}
For $0 < p < 1$, the $p$-th \emph{quantile} of $X$ is the value $x_p$ with
$F(x_p) = p$, i.e.
\[
  x_p = F^{-1}(p),
\]
when $F$ is continuous and strictly increasing so the inverse exists. The $0.5$
quantile is the \emph{median}; the quartiles are $p = 0.25, 0.5, 0.75$; percentiles
use $p = k/100$.
\end{definition}

Geometrically, the median splits the area under the density into two equal halves of
$0.5$ each. The quartiles cut it into four equal-area quarters. Quantiles are the
``area coordinates'' of the distribution. Figure~\ref{ch10:fig-median} shows the
median as the vertical line that balances area, not the line that balances width.

\begin{figure}[H]
\centering
\begin{tikzpicture}
\begin{axis}[
  width=0.85\textwidth, height=5.4cm,
  title style={font=\small},
  title={Median Splits Area Under the Density Curve into Equal Halves},
  xlabel={$x$}, ylabel={$f(x)=3x^2$},
  xmin=-0.05, xmax=1.1, ymin=0, ymax=3.3,
  domain=0:1, samples=100,
  axis lines=left, ytick=\empty,
  xtick={0.7937}, xticklabels={$m\approx0.794$},
]
\addplot[name path=curve, densityred, very thick] {3*x^2};
\path[name path=axisL] (axis cs:0,0) -- (axis cs:0.7937,0);
\addplot[densityred!22] fill between[of=curve and axisL, soft clip={domain=0:0.7937}];
\addplot[black, dashed, thick] coordinates {(0.7937,0)(0.7937,1.8907)};
\node at (axis cs:0.55,0.5) {\small area $=0.5$};
\node at (axis cs:0.9,1.9) {\small area $=0.5$};
\end{axis}
\end{tikzpicture}
\caption{The median $m$ of $f(x)=3x^2$ on $[0,1]$ splits the area into two equal
halves of $0.5$. Because the density piles up on the right, the equal-area line sits
well to the right of the midpoint $0.5$ of the interval.}
\label{ch10:fig-median}
\end{figure}

\begin{example}[Quantiles of $f=3x^2$]
\label{ch10:ex-quantile}
We found $F(x) = x^3$ on $[0,1]$. Invert it: $F(x)=p$ means $x^3 = p$, so
\[
  x_p = p^{1/3}.
\]
The median is $x_{0.5} = 0.5^{1/3} \approx 0.794$ -- notably larger than the mean
$\E[X]=0.75$ from Example~\ref{ch10:ex-EV}. \emph{Mean below median} is the signature
of a left-skewed shape: here the density rises toward $1$ and has a long thin tail
toward $0$, which drags the mean down more than the median. The lower and upper
quartiles are $0.25^{1/3}\approx 0.630$ and $0.75^{1/3}\approx 0.909$.
\end{example}

\begin{trapbox}
The mean and the median are different summaries and can disagree a lot. The mean is
the balance point and is sensitive to far-away tail mass; the median is the
equal-area split and shrugs off the tails. For symmetric densities they coincide.
For skewed ones -- incomes, waiting times, file sizes -- they can be far apart, which
is exactly why ``median household income'' is reported instead of the mean: a few
billionaires move the mean but barely budge the median.
\end{trapbox}

The quantile function is also a practical tool for generating random variables,
allowing computer algorithms to produce samples from any continuous distribution
starting from standard uniform draws.

\begin{theorem}[Inverse-transform / probability integral transform]
\label{ch10:thm-pit}
Let $F$ be a continuous, strictly increasing CDF with inverse $F^{-1}$.
\begin{enumerate}
\item \textbf{(Generation.)} If $U \sim \Unif(0,1)$, then $X = F^{-1}(U)$ has CDF
      $F$.
\item \textbf{(Flattening.)} Conversely, if $X$ has CDF $F$, then $F(X) \sim
      \Unif(0,1)$.
\end{enumerate}
\end{theorem}

\begin{whybox}
Both halves are one short computation. For generation, compute the CDF of
$X = F^{-1}(U)$ directly, using that $F$ is increasing (so it preserves the
inequality) and that $U$ is uniform (so $\P(U \le u) = u$):
\[
  \P(X \le x) = \P\!\big(F^{-1}(U) \le x\big) = \P\big(U \le F(x)\big) = F(x).
\]
That is exactly the CDF we wanted, so $X \sim F$. For flattening, run it backward:
\[
  \P\big(F(X) \le p\big) = \P\big(X \le F^{-1}(p)\big) = F\big(F^{-1}(p)\big) = p,
\]
which is the CDF of $\Unif(0,1)$. The deep content is that \emph{a CDF maps its own
distribution to the flat uniform}: pushing $X$ through its accumulated-probability
function evens out all the lumps, because $F(x)$ measures exactly how much
probability sits below $x$. Generation simply reverses that flattening.

Figure~\ref{ch10:fig-pit} visualizes the mechanism. Draw a uniform height
$U \in [0,1]$ on the vertical axis, shoot horizontally to the CDF ramp $y=F(x)$,
and drop straight down to read $X = F^{-1}(U)$ on the horizontal axis. Where the
CDF is steep (tall density $f = F'$), a wide vertical slice of $U$-values maps into
a narrow horizontal window, densely clustering samples. Where the CDF is nearly
flat (low density), a narrow vertical slice is stretched across a wide range of
$x$, leaving samples sparse. The ramp's slope automatically generates the correct
density.
\end{whybox}

\begin{figure}[H]
\centering
\begin{tikzpicture}
\begin{axis}[
  width=0.85\textwidth, height=5.8cm,
  title style={font=\small},
  title={Inverse-Transform Sampling: Generating $X = F^{-1}(U)$ via the CDF Ramp},
  xlabel={$x$}, ylabel={$u = F(x)$},
  xmin=-0.05, xmax=1.12, ymin=-0.05, ymax=1.08,
  domain=0:1, samples=100,
  axis lines=left,
  ytick={0, 0.125, 0.5, 1}, yticklabels={$0$, $0.125$, $0.5$, $1$},
  xtick={0, 0.5, 0.7937, 1}, xticklabels={$0$, $0.5$, $x_m \approx 0.794$, $1$},
]
\addplot[densityred, very thick] {x^3};
\addplot[black, dashed, thick, -{Stealth}] coordinates {(0,0.5)(0.7937,0.5)};
\addplot[black, dashed, thick, -{Stealth}] coordinates {(0.7937,0.5)(0.7937,0)};
\addplot[quotientorange, dashed, thick, -{Stealth}] coordinates {(0,0.125)(0.5,0.125)};
\addplot[quotientorange, dashed, thick, -{Stealth}] coordinates {(0.5,0.125)(0.5,0)};
\node[above, font=\footnotesize] at (axis cs:0.38, 0.51) {\textbf{1.} Project $u=0.5$ across};
\node[right, font=\footnotesize] at (axis cs:0.80, 0.25) {\textbf{2.} Drop to $x_m$};
\node[draw=gray!50, fill=white, rounded corners=2pt, font=\scriptsize, align=left] at (axis cs:0.36, 0.84)
  {\textbf{Flat ramp:} sparse samples ($12.5\%$ in $[0,0.5]$)\\[2pt]\textbf{Steep ramp:} dense samples ($87.5\%$ in $[0.5,1]$)};
\end{axis}
\end{tikzpicture}
\caption{The geometry of inverse-transform sampling. Pick $U \sim \Unif(0,1)$ on the
vertical axis, project across to the curve $y = F(x)$, and drop down to the $x$-axis
to find $X = F^{-1}(U)$. For $F(x)=x^3$, the bottom half $[0, 0.5]$ of the $x$-domain
receives only the lowest $1/8$ ($12.5\%$) of the uniform height, while the upper half
$[0.5, 1]$ catches the remaining $7/8$ ($87.5\%$), clustering samples where the density
$f(x)=3x^2$ is tall.}
\label{ch10:fig-pit}
\end{figure}

\begin{applicationbox}
\textbf{Simulating any distribution from coin-flip randomness.} A computer's
random-number generator hands you $\Unif(0,1)$ samples and nothing else. The
inverse-transform theorem turns that single source into samples from \emph{any}
distribution whose quantile function you can write down. Want an exponential waiting
time with rate $\lambda$? Its CDF is $F(x)=1-e^{-\lambda x}$, so
$F^{-1}(u) = -\tfrac{1}{\lambda}\ln(1-u)$. Draw $U\sim\Unif(0,1)$ and report
$X = -\tfrac{1}{\lambda}\ln(1-U)$; the result is exactly $\Expo(\lambda)$. This principle
forms the foundation of Monte Carlo simulation, which we develop further in the
chapter on Monte Carlo methods.
\end{applicationbox}

\section{A cautionary paradox}
\label{ch10:paradox}

Continuous probability rewards careful thinking and punishes sloppy phrases like
``choose a random chord.'' Here is the classic warning.

\begin{paradoxbox}
\textbf{Bertrand's chord paradox (1889).} Draw a circle and inscribe an equilateral
triangle. Pick a chord ``at random.'' What is the probability the chord is longer
than a side of the triangle?

Three perfectly reasonable people give three different answers:
\begin{itemize}
\item \emph{Random endpoints.} Fix one endpoint, choose the other uniformly on the
      circle. The chord beats the side when the second point lands on the far arc:
      probability $\tfrac{1}{3}$.
\item \emph{Random radius.} Choose a radius, then a point uniformly along it for the
      chord's midpoint. The chord is long enough when the midpoint is within half a
      radius of the center: probability $\tfrac{1}{2}$.
\item \emph{Random midpoint.} Choose the chord's midpoint uniformly over the disk's
      \emph{area}. The chord is long when the midpoint falls inside the inscribed
      circle of half the radius, whose area is one-quarter of the disk: probability
      $\tfrac{1}{4}$.
\end{itemize}
All three are correct -- for the procedure they describe. \emph{Resolution:} ``at
random'' is meaningless on a continuum until you specify \emph{which density} you are
sampling from. The three answers correspond to three genuinely different
distributions on the space of chords. In the discrete world ``equally likely'' has
one obvious meaning (count and divide); in the continuous world there is no single
canonical uniform measure on a curved or higher-dimensional space, and you must name
your density before any probability is even defined.
\end{paradoxbox}

\begin{happeningbox}
The moral generalizes far beyond circles. A continuous probability question is
\emph{not well posed} until the density (equivalently, the CDF) is on the table.
``Pick a random point,'' ``a random line,'' ``a random triangle'' -- each hides a
choice of measure. Whenever you meet ``at random'' in a continuous setting, your
first job is to ask: \emph{random according to what density?} Naming $f$ is not a
technicality; it is what makes the probability exist at all.
\end{happeningbox}

\section{Mixed and degenerate cases (a brief caveat)}
\label{ch10:mixed}

Not every random variable is purely discrete or purely continuous. Consider an
insurance payout that is $0$ with probability $0.7$ (no claim filed) and otherwise a
continuous positive amount. Its CDF jumps at $0$ (a genuine point mass of $0.7$) and
then rises smoothly -- it is a \emph{mixture}. Such variables have no ordinary
density, because no integrable curve can produce a jump.

\begin{marginnotebox}
This is exactly why the CDF, not the pdf, is the truly universal object: \emph{every}
random variable has a CDF, but only purely continuous ones have a density and only
discrete ones have a PMF. The CDF handles staircases, ramps, and any blend of the two
in one uniform language. When we build measure-theoretic probability later, the CDF
(and the measure it induces) becomes the bedrock, with PMFs and pdfs as two special
cases distinguished by whether the CDF is a sum of jumps or a smooth integral. For
the rest of this book, ``continuous'' will mean ``has a density'' unless we say
otherwise.
\end{marginnotebox}

\section{Chapter summary}
\label{ch10:summary}

\begin{focusbox}
\textbf{The continuous toolkit.}
\begin{itemize}
\item A continuous $X$ has a \textbf{density} $f \ge 0$ with $\int_{-\infty}^\infty
      f = 1$, and probability is \textbf{area}: $\P(a\le X\le b)=\int_a^b f$.
\item \textbf{Points carry no probability:} $\P(X=x)=0$, so $<$ and $\le$ are
      interchangeable; $f(x)$ is probability \emph{per unit length}, not a
      probability, and may exceed $1$.
\item The \textbf{CDF} $F(x)=\P(X\le x)=\int_{-\infty}^x f$ is universal; it is a
      jumpless ramp for continuous $X$. Recover the density by differentiating:
      $f = F'$ (by the Fundamental Theorem of Calculus).
\item \textbf{Expectation and variance are integrals:} $\E[X]=\int x f$,
      $\E[g(X)]=\int g\,f$, $\Var(X)=\int(x-\mu)^2 f = \E[X^2]-\mu^2$. All the old
      identities survive, with $\sum \to \int$.
\item The \textbf{uniform} $\Unif(a,b)$ is the base case: flat density
      $1/(b-a)$, mean $(a+b)/2$, variance $(b-a)^2/12$.
\item \textbf{Quantiles} invert the CDF, $x_p=F^{-1}(p)$; the median is the
      equal-area split. The inverse-transform $X=F^{-1}(U)$ turns uniform noise into
      any distribution.
\item ``At random'' in the continuum is \textbf{undefined until you name the
      density} (Bertrand's paradox).
\end{itemize}
\end{focusbox}

\begin{connectionbox}
We can now model continuous measurements and translate questions about chance into
calculus. In practice, many systems cluster into recurring families of distributions:
memoryless waiting times, heavy-tailed extremes, and bounded fractions.

Chapter~11 introduces these parametric families -- the Exponential, Gamma, Beta, and
Cauchy distributions -- before Chapter~12 explores the Normal distribution.
\end{connectionbox}

\section{Exercises}
\label{ch10:exercises}

\begin{exercise}[Normalize and use a density]
\label{ch10:exr-norm}
Let $f(x) = c(1-x^2)$ for $-1 \le x \le 1$ and $0$ otherwise.
(a) Find $c$. (b) Compute $\P(X > 1/2)$. (c) By symmetry, what is $\E[X]$? Verify
with the integral.
\end{exercise}

\begin{exercise}[CDF to pdf and back]
\label{ch10:exr-cdf}
A variable has CDF $F(x) = \sqrt{x}$ for $0 \le x \le 1$ (with $0$ below, $1$ above).
(a) Find the density $f$. (b) Sketch why $f$ blows up near $0$ even though it is a
valid density. (c) Compute $\P(0.25 \le X \le 0.81)$ two ways: via $F$, and via
$\int f$.
\end{exercise}

\begin{exercise}[Exponential basics]
\label{ch10:exr-expo}
Let $X$ have density $f(x) = \lambda e^{-\lambda x}$ for $x \ge 0$. (a) Find the CDF.
(b) Show $\E[X] = 1/\lambda$ (integrate by parts). (c) Show $\Var(X)=1/\lambda^2$.
(d) Find the median and explain why it is smaller than the mean.
\end{exercise}

\begin{exercise}[Why $\P(X=x)=0$ in your own words]
\label{ch10:exr-pointmass}
Using the area picture, explain to a skeptical friend how every outcome of a
continuous variable can have probability $0$ while one of them certainly occurs.
Then give the $\eps$-trapping argument from Section~\ref{ch10:point-mass-zero} in two
or three lines.
\end{exercise}

\begin{exercise}[Uniform computations]
\label{ch10:exr-unif}
Let $X \sim \Unif(2,8)$. (a) Write $f$ and $F$. (b) Find $\P(3 \le X \le 5)$.
(c) Find $\E[X]$ and $\Var(X)$ from the formulas, then re-derive $\Var$ by direct
integration as a check. (d) Find the $0.9$ quantile.
\end{exercise}

\begin{exercise}[Mean versus median under skew]
\label{ch10:exr-skew}
For $f(x)=4x^3$ on $[0,1]$: (a) find the CDF; (b) find the mean and the median;
(c) state which is larger and connect your answer to the direction of the skew, as
in Example~\ref{ch10:ex-quantile}.
\end{exercise}

\begin{exercise}[Simulating by inverse transform]
\label{ch10:exr-pit}
You can only draw $U \sim \Unif(0,1)$. (a) Give an explicit formula $X=g(U)$ that
produces $X \sim \Unif(2,8)$. (b) Give a formula that produces a variable with CDF
$F(x)=x^4$ on $[0,1]$. (c) Verify (b) by computing the CDF of your $g(U)$ directly,
using $\P(U\le u)=u$.
\end{exercise}

\begin{exercise}[Density can exceed one]
\label{ch10:exr-bignum}
Let $X \sim \Unif(0, 0.01)$. (a) What is the height of its density? (b) Reconcile a
density value of $100$ with the rule that probabilities never exceed $1$.
(c) Compute $\P(X \le 0.001)$ and confirm it is a legitimate probability.
\end{exercise}

\begin{exercise}[Proof: variance shortcut]
\label{ch10:exr-varproof}
Prove for a continuous $X$ with mean $\mu$ that
$\int_{-\infty}^{\infty}(x-\mu)^2 f(x)\,\dd x = \int_{-\infty}^{\infty} x^2 f(x)\,\dd x
- \mu^2$. State exactly which two properties of the integral you use, and where the
fact $\int f = 1$ enters.
\end{exercise}

\begin{exercise}[Conceptual: Bertrand revisited]
\label{ch10:exr-bertrand}
In your own words, explain why Bertrand's chord paradox is \emph{not} a flaw in
probability theory but a flaw in the problem statement. What single piece of
information, if supplied, would make the question have exactly one correct answer?
\end{exercise}
